% Part: lambda-calculus
% Chapter: introduction
% Section: overview

\documentclass[../../../include/open-logic-section]{subfiles}

\begin{document}

\olfileid{lam}{int}{ovr}
\olsection{Gambaran Umum}

Kalkulus lambda wiwitane dirancang dening Alonzo Church
ing wiwitan dasawarsa 1930-an minangka dhasar logika
konstruktif, lan \emph{dudu} minangka model fungsi
kang bisa diitung. Nanging ora suwe banjur dibuktekake
yen kalkulus iki ekuivalen karo definisi komputabilitas
liyane, kayata fungsi kang bisa diitung nganggo mesin
Turing lan fungsi rekursif parsial. Kasunyatan yen
asil iki wiwitane rada nggumunake ndadekake
karakterisasi kasebut luwih narik kawigaten.

Notasi lambda iku cara kang praktis kanggo nyebut
fungsi langsung nganggo ekspresi simbolis kang
netepake fungsi iku, tanpa kudu luwih dhisik
menehi jeneng kapisah. Tinimbang nyebut
``upamana $f$ iku fungsi kang ditetepake
dening $f(x) = x + 3$,'' kita bisa nyebut
``upamana $f$ iku fungsi
$\lambd[x][(x + 3)]$.'' Kanthi tembung liya,
$\lambd[x][(x+3)]$ mung sawijining \emph{jeneng}
kanggo fungsi kang nambah telu marang argumene.
Ing ekspresi iki, $x$ iku variabel semu,
utawa panggonan isian: fungsi kang padha
uga bisa ditulis minangka $\lambd[y][(y + 3)]$.
Notasi iki uga lumaku yen ana parameter liyane.
Umpamane, upamana $g(x, y)$ iku fungsi rong
variabel, lan $k$ iku wilangan asli. Banjur
$\lambd[x][g(x,k)]$ iku fungsi kang
ngirim saben~$x$ menyang~$g(x, k)$.

Cara netepake fungsi saka ekspresi simbolis iki
diarani \emph{abstraksi lambda}. Kosok baline
abstraksi lambda yaiku \emph{aplikasi}:
yen kita duwe fungsi $f$ (umpamane, kang
ditetepake ing wilangan asli), kita bisa
\emph{ngaplikasikake} fungsi iku marang
sembarang nilai, kayata 2. Ing notasi
biasa, asile mesthi ditulis~$f(2)$.

Apa kang kedadeyan yen abstraksi lambda
digabung karo aplikasi? Ekspresi kang
diasilake banjur bisa disederhanakake
kanthi ``ngiseni'' variabel kang diabstraksi
nganggo argumen aplikasi. Umpamane,
\[
(\lambd[x][(x + 3)])(2)
\]
bisa disederhanakake dadi~$2 + 3$.

Nganti kene kita mung ngenalake notasi anyar
kanggo gagasan kang wis lumrah. Nanging kalkulus
lambda luwih adoh ninggalake sudut pandang
teori himpunan. Ing kerangka iki:
\begin{enumerate}
\item Saben apa wae nyebut sawijining fungsi.
\item Fungsi bisa ditetepake nganggo abstraksi lambda.
\item Apa wae bisa diaplikasikake marang apa wae liyane.
\end{enumerate}
Umpamane, yen $F$ iku suku ing kalkulus lambda,
$F(F)$ tansah dianggep nduweni teges.
Kerangka kang bebas iki diarani kalkulus lambda
\emph{tanpa tipe}, ing kene ``tanpa tipe''
ateges ``ora ana watesan ngenani apa kang
bisa diaplikasikake marang apa''.

\begin{digress}
Ana uga kalkulus lambda \emph{mawa tipe},
yaiku variasi wigati saka versi tanpa tipe.
Sanadyan kalkulus lambda mawa tipe padha
karo versi tanpa tipe ing akeh bab, kalkulus
mawa tipe luwih gampang diselarasake karo
kerangka teori himpunan klasik lan nduweni
sawetara sipat kang beda banget.

Panliten ngenani kalkulus lambda wis dadi
bagean utama ing ilmu komputer teoretis
lan ing rancangan basa pamrograman. LISP,
kang dirancang dening John McCarthy ing
dasawarsa 1950-an, minangka tuladha wiwitan
basa kang kena pengaruh saka gagasan-gagasan iki.
\end{digress}

\end{document}
